% !TEX root = ../diff_equations.tex
% =================================================================
%  sections/lecture03.tex
%  Лекция 3. Лемма Гронуолла. Теорема Пикара. Метод сжимающих отображений.
% =================================================================

\part*{Лекция 3}

\section{Лемма Гронуолла. Теорема Пикара. Метод сжимающих отображений}

\subsection{Лемма Гронуолла}\label{sec:гронуолл}

\begin{Lemma}{Гронуолла}{гронуолл}
	Пусть $\vphi(t) \geqslant 0$, $\vphi \in C(a,b)$, $t_0 \in (a,b)$ и
	$\lambda, \mu \geqslant 0$, причём
	\[
	\vphi(t) \leqslant \lambda + \mu\abs{\int_{t_0}^{t} \vphi(s)\,\dd s},
	\qquad t \in (a,b).
	\]
	Тогда
	\[
	\vphi(t) \leqslant \lambda e^{\mu\abs{t - t_0}}, \qquad t \in (a,b).
	\]
\end{Lemma}

\begin{proof}
	Разберём случай $t \geqslant t_0$. Положим
	\[
	\Phi(t) = \lambda + \mu\int_{t_0}^{t} \vphi(s)\,\dd s,
	\qquad\text{так что}\qquad \vphi(t) \leqslant \Phi(t).
	\]
	Тогда $\dot\Phi(t) = \mu\vphi(t) \leqslant \mu\Phi(t)$, откуда
	\[
	e^{-\mu(t - t_0)}\dot\Phi - \mu e^{-\mu(t-t_0)}\Phi \leqslant 0,
	\qquad
	\dder{}{t}\Bigl(e^{-\mu(t-t_0)}\Phi\Bigr) \leqslant 0.
	\]
	Значит, функция $e^{-\mu(t-t_0)}\Phi(t)$ не возрастает, поэтому
	\[
	e^{-\mu(t-t_0)}\Phi(t) \leqslant \Phi(t_0) = \lambda,
	\qquad
	\vphi(t) \leqslant \Phi(t) \leqslant \lambda e^{\mu(t - t_0)}.
	\]

	\medskip
	Случай $t \leqslant t_0$ разбирается аналогично.
\end{proof}

% TODO (сверить с лектором): случай $t \leqslant t_0$ в лемме Гронуолла был
% оставлен слушателям; в конспекте он не разобран.

% ---------------------------------------------------------------
% РИСУНОК 1: «воронка» Гронуолла — область под барьером
% lambda e^{mu |t - t_0|} (lambda = 0.7, mu = 1) и функция phi внутри неё.
% ---------------------------------------------------------------
\begin{figure}[H]
	\centering
	\begin{tikzpicture}[x=2.0cm, y=0.85cm, >=Stealth]
		% допустимая область: под барьером
		\fill[hlBlue, hlarea] (-1.5,0)
		-- plot[domain=-1.5:1.5, samples=140] (\x,{0.7*exp(abs(\x))})
		-- (1.5,0) -- cycle;

		\draw[ax] (-1.75,0) -- (1.9,0) node[below right] {$t$};
		\draw[ax] (0,-0.4) -- (0,3.75) node[above] {$\vphi$};

		% барьер и сама функция
		\draw[curvehl, domain=-1.5:1.5, samples=140]
		plot (\x,{0.7*exp(abs(\x))});
		\draw[curve, domain=-1.5:1.5, samples=100]
		plot (\x,{0.7*(1 + 0.45*\x*\x)});

		\draw[tickmark] (0.05,0.7) -- (-0.05,0.7);
		\node[slbl, left]  at (-0.07,0.7) {$\lambda$};
		\node[slbl, below] at (0,-0.12) {$t_0$};
		\node[pt] at (0,0.7) {};

		\node[Purple, slbl, right] at (1.55,3.15) {$\lambda e^{\mu\abs{t-t_0}}$};
		\node[Blue, slbl, right]   at (1.55,1.42) {$\vphi(t)$};
	\end{tikzpicture}
	\caption{Лемма Гронуолла загоняет $\vphi$ в «воронку»: интегральное
		неравенство означает, что график не может выйти из голубой области под
		барьером $\lambda e^{\mu\abs{t-t_0}}$ (выделенная кривая). В точке $t_0$
		воронка сужается до значения $\lambda$, а с удалением от неё расширяется
		экспоненциально.}
	\label{fig:гронуолл}
\end{figure}

\begin{Remark}[addon]{Оценка неулучшаема}{гронуолл-точность}
	Барьер на рис.~\ref{fig:гронуолл} нельзя опустить: на функции
	$\vphi(t) = \lambda e^{\mu(t-t_0)}$ исходное неравенство обращается в
	равенство, потому что
	\[
	\lambda + \mu\int_{t_0}^{t} \lambda e^{\mu(s-t_0)}\,\dd s
	= \lambda + \lambda\bigl(e^{\mu(t-t_0)} - 1\bigr)
	= \lambda e^{\mu(t-t_0)} = \vphi(t).
	\]
	То есть заключение леммы достигается, и множитель $e^{\mu\abs{t-t_0}}$
	заменить на меньший нельзя.
\end{Remark}

\begin{Corollary}{Нулевая оценка}{гронуолл-следствие}
	При $\lambda = 0$ из
	\[
	0 \leqslant \vphi(t) \leqslant \mu\abs{\int_{t_0}^{t} \vphi(s)\,\dd s}
	\]
	следует $\vphi(t) \equiv 0$.
\end{Corollary}

\subsection{Теорема Пикара}\label{sec:пикар}

\begin{Theorem}{Пикара}{пикар}
	Пусть $\dot x = f(t,x)$, где $f \in C(G)$ и
	$f \in \mathrm{Lip}_{x,\mathrm{loc}}(G)$. Тогда $G$ --- область
	существования и единственности.
\end{Theorem}

\subsubsection*{Существование: последовательные приближения Пикара}

Пусть $\alpha, \beta > 0$ таковы, что
\[
R = \bigl\{(t,x) \colon \abs{t - t_0} \leqslant \alpha,\
\abs{x - x_0} \leqslant \beta\bigr\} \subset G,
\]
и пусть $\abs{f(t,x)} \leqslant M$ в $R$, $h = \min\bigl(\alpha, \beta/M\bigr)$ ---
то же построение, что и на рис.~\ref{fig:промежуток-пеано}, с промежутком
Пеано $[t_0-h,\ t_0+h]$.

\textbf{Последовательные приближения Пикара:}
\[
\vphi_0(t) \equiv x_0, \qquad t \in [t_0-h,\ t_0+h],
\]
\[
\vphi_{k+1}(t) = x_0 + \int_{t_0}^{t} f\bigl(s, \vphi_k(s)\bigr)\,\dd s,
\qquad k \geqslant 0.
\]

\begin{Lemma}{Приближения определены и не выходят из $R$}{приближения-определены}
	Все функции $\vphi_k(t)$ определены на $[t_0-h,\ t_0+h]$, причём
	\[
	(t, \vphi_k(t)) \in R, \qquad t \in [t_0-h,\ t_0+h].
	\]
\end{Lemma}

\begin{proof}
	При $k = 0$ утверждение очевидно. Переход $k \mapsto k+1$: из
	$(t, \vphi_k(t)) \in R \subset G$ следует, что $f\bigl(t, \vphi_k(t)\bigr)$
	определена при $t \in [t_0-h,\ t_0+h]$, и тогда
	\[
	\abs{\vphi_{k+1}(t) - x_0}
	\leqslant \abs{\int_{t_0}^{t} f\bigl(s, \vphi_k(s)\bigr)\,\dd s}
	\leqslant M\abs{t - t_0} \leqslant Mh \leqslant \beta.
	\]
\end{proof}

\begin{Example}[addon]{Приближения Пикара для $\dot x = x$}{приближения-пример}
	Пусть $\dot x = x$, $x(0) = 1$. Тогда $\vphi_0 \equiv 1$ и
	\[
	\vphi_{k+1}(t) = 1 + \int_0^t \vphi_k(s)\,\dd s ,
	\]
	откуда
	\[
	\vphi_1 = 1 + t, \qquad
	\vphi_2 = 1 + t + \frac{t^2}{2}, \qquad
	\vphi_3 = 1 + t + \frac{t^2}{2} + \frac{t^3}{6} .
	\]

	\medskip
	Индукцией получаем $\vphi_k(t) = \sum_{j=0}^{k} \dfrac{t^j}{j!}$ --- частичные
	суммы ряда для $e^t$. Приближения Пикара здесь в точности воспроизводят
	разложение точного решения $x(t) = e^t$ (рис.~\ref{fig:приближения-пикара}).
\end{Example}

% ---------------------------------------------------------------
% РИСУНОК 2: приближения Пикара phi_0, ..., phi_3 для x' = x, x(0) = 1
% на [0; 2] и точное решение e^t.
% ---------------------------------------------------------------
\begin{figure}[H]
	\centering
	\begin{tikzpicture}[x=2.0cm, y=0.4cm, >=Stealth]
		\draw[ax] (-0.15,0) -- (2.5,0) node[below right] {$t$};
		\draw[ax] (0,-0.6) -- (0,8.4) node[above] {$x$};

		% точное решение
		\draw[curvehl, domain=0:2, samples=120] plot (\x,{exp(\x)});

		% приближения Пикара
		\draw[curve] (0,1) -- (2,1);
		\draw[curve] (0,1) -- (2,3);
		\draw[curve, domain=0:2, samples=80]
		plot (\x,{1 + \x + 0.5*\x*\x});
		\draw[curve, domain=0:2, samples=80]
		plot (\x,{1 + \x + 0.5*\x*\x + \x*\x*\x/6});

		\node[pt] at (0,1) {};
		\draw[tickmark] (0.04,1) -- (-0.04,1);
		\node[slbl, left] at (-0.06,1) {$1$};
		\node[slbl, below] at (2,-0.15) {$2$};

		\node[Purple, slbl, right] at (2.06,7.389) {$e^t$};
		\node[Blue, slbl, right]   at (2.06,6.333) {$\vphi_3$};
		\node[Blue, slbl, right]   at (2.06,5.0)   {$\vphi_2$};
		\node[Blue, slbl, right]   at (2.06,3.0)   {$\vphi_1$};
		\node[Blue, slbl, right]   at (2.06,1.0)   {$\vphi_0$};
	\end{tikzpicture}
	\caption{Приближения Пикара для $\dot x = x$, $x(0) = 1$: каждое следующее
		добавляет один член ряда, и $\vphi_k$ --- это частичная сумма для $e^t$
		(выделенная кривая). С ростом $k$ приближения поднимаются к точному
		решению, причём быстрее всего вблизи начальной точки.}
	\label{fig:приближения-пикара}
\end{figure}

\subsubsection*{Сходимость последовательных приближений}

\begin{Lemma}{Сходимость приближений Пикара}{сходимость-приближений}
	Приближения $\vphi_k(t)$ равномерно сходятся на $[t_0-h,\ t_0+h]$.
\end{Lemma}

\begin{proof}
	\textbf{Шаг 1 (переход к ряду).} Положим
	\[
	\psi_0(t) = \vphi_0(t), \qquad
	\psi_k(t) = \vphi_k(t) - \vphi_{k-1}(t), \quad k \geqslant 1,
	\]
	и рассмотрим ряд
	\begin{equation}\label{eq:ряд-пикара}
		\sum_{k=0}^{\infty} \psi_k(t).
	\end{equation}
	Его частичная сумма телескопируется:
	\[
	\Sigma_m(t) = \sum_{k=0}^{m} \psi_k(t)
	= \vphi_0 + (\vphi_1 - \vphi_0) + \dots + (\vphi_m - \vphi_{m-1})
	= \vphi_m(t),
	\]
	поэтому сходимость последовательности $\{\vphi_k\}$ равносильна сходимости
	ряда~\eqref{eq:ряд-пикара}.

	\medskip
	\textbf{Шаг 2 (оценка членов ряда).} Пусть $L > 0$ --- константа Липшица
	функции $f$ по $x$ на $R$ (она существует по лемме об условии Липшица на
	компакте). Докажем, что для $k \geqslant 1$ и $t \in [t_0-h,\ t_0+h]$
	\begin{equation}\label{eq:оценка-psi}
		\abs{\psi_k(t)} \leqslant \frac{M}{L}\,\frac{L^k\abs{t-t_0}^k}{k!} .
	\end{equation}
	Разберём случай $t \in [t_0,\ t_0+h]$; при $t \leqslant t_0$ рассуждение
	такое же. Индукция по $k$.

	\medskip
	\textbf{База $k = 1$.} Так как $\vphi_0 \equiv x_0$,
	\[
	\abs{\psi_1(t)} = \abs{\int_{t_0}^{t} f(s, x_0)\,\dd s}
	\leqslant M\abs{t - t_0} = \frac{M}{L}\,\frac{L\abs{t-t_0}}{1!} .
	\]

	\medskip
	\textbf{Переход $k \mapsto k+1$.} По определению приближений и условию
	Липшица
	\[
	\abs{\psi_{k+1}(t)}
	= \abs{\int_{t_0}^{t} \bigl[f\bigl(s, \vphi_k(s)\bigr)
	- f\bigl(s, \vphi_{k-1}(s)\bigr)\bigr]\,\dd s}
	\leqslant L\int_{t_0}^{t} \abs{\vphi_k(s) - \vphi_{k-1}(s)}\,\dd s ,
	\]
	а подынтегральное выражение --- это $\abs{\psi_k(s)}$. Подставляя
	предположение индукции~\eqref{eq:оценка-psi}, получаем
	\[
	\abs{\psi_{k+1}(t)}
	\leqslant L\int_{t_0}^{t} \frac{M}{L}\,\frac{L^k(s-t_0)^k}{k!}\,\dd s
	= \frac{M}{L}\,L^{k+1}\,\frac{(t-t_0)^{k+1}}{(k+1)!} .
	\]

	\medskip
	\textbf{Шаг 3 (мажорирование).} На всём промежутке
	$\abs{t - t_0} \leqslant h$, поэтому
	\[
	\abs{\psi_k(t)} \leqslant A_k := \frac{M}{L}\,\frac{(Lh)^k}{k!} ,
	\]
	а числовой ряд $\sum_k A_k$ сходится: его сумма не превосходит
	$\dfrac{M}{L}e^{Lh}$. По признаку Вейерштрасса
	ряд~\eqref{eq:ряд-пикара} сходится равномерно, то есть
	$\vphi_k(t) \rightrightarrows \vphi(t)$ на $[t_0-h,\ t_0+h]$.

	\medskip
	\textbf{Шаг 4 (предельный переход).} Прямоугольник $R$ --- компакт,
	поэтому $f$ на нём равномерно непрерывна, и из $\vphi_k \rightrightarrows \vphi$
	следует $f\bigl(t, \vphi_k(t)\bigr) \rightrightarrows f\bigl(t, \vphi(t)\bigr)$.
	Переходя к пределу в равенстве
	\[
	\vphi_{k+1}(t) = x_0 + \int_{t_0}^{t} f\bigl(s, \vphi_k(s)\bigr)\,\dd s ,
	\]
	получаем
	\[
	\vphi(t) = x_0 + \int_{t_0}^{t} f\bigl(s, \vphi(s)\bigr)\,\dd s ,
	\]
	то есть $\vphi$ --- решение интегрального уравнения~\eqref{eq:эиу}, а значит
	и решение задачи Коши.
\end{proof}

\subsubsection*{Единственность}

Пусть $(t_0,x_0) \in G$, прямоугольник $R$ --- тот же, что и выше, а $L$ ---
константа Липшица $f$ по $x$ в $R$. Пусть $x_1(t)$ и $x_2(t)$ --- решения
задачи Коши с начальными данными $(t_0,x_0)$. По непрерывности найдётся
$(a,b) \ni t_0$, на котором графики обоих решений лежат в $R$:
\[
\bigl(t, x_1(t)\bigr),\ \bigl(t, x_2(t)\bigr) \in R,
\qquad t \in (a,b).
\]

Оба решения удовлетворяют интегральному уравнению:
\[
x_i(t) = x_0 + \int_{t_0}^{t} f\bigl(s, x_i(s)\bigr)\,\dd s,
\qquad i = 1,2 .
\]

Положим $\vphi(t) = \abs{x_1(t) - x_2(t)}$; эта функция неотрицательна и
непрерывна, и
\[
\vphi(t)
= \abs{\int_{t_0}^{t} \bigl[f\bigl(s, x_1(s)\bigr)
- f\bigl(s, x_2(s)\bigr)\bigr]\,\dd s}
\leqslant \abs{\int_{t_0}^{t} \abs{f\bigl(s, x_1(s)\bigr)
- f\bigl(s, x_2(s)\bigr)}\,\dd s}
\leqslant L\abs{\int_{t_0}^{t} \vphi(s)\,\dd s} .
\]

Это в точности посылка следствия~\ref{cor:гронуолл-следствие} из леммы
Гронуолла при $\lambda = 0$, $\mu = L$, поэтому $\vphi(t) \equiv 0$, то есть
$x_1(t) \equiv x_2(t)$.

\begin{Remark}[addon]{Пеано и Пикар: что даёт каждая}{пеано-пикар}
	\textbf{Условие.} Пеано требует только $f \in C(G)$; Пикар добавляет к
	этому локальное условие Липшица по $x$.

	\medskip
	\textbf{Вывод.} Пеано --- существование решения задачи Коши; Пикар ---
	существование и единственность, то есть $G$ оказывается областью
	существования и единственности.

	\medskip
	\textbf{Метод.} У Пеано --- ломаные Эйлера и компактность, причём сходится
	лишь подпоследовательность. У Пикара --- последовательные приближения,
	сходящиеся все сразу, а единственность вытекает из леммы Гронуолла.

	\medskip
	\textbf{Где видна разница.} У уравнения $\dot x = 3x^{2/3}$ правая часть
	непрерывна, но не липшицева в нуле: решение через начало координат
	существует по Пеано, однако оно не одно (лекция~1). Условие Липшица такую
	картину исключает.
\end{Remark}

\subsection{Глобальная теорема единственности}\label{sec:глобальная-единственность}

\begin{Theorem}{Глобальная единственность}{глобальная-единственность}
	Пусть $G$ --- область существования и единственности, а $x_1(t)$ и $x_2(t)$
	--- решения на $(a,b)$, совпадающие в одной точке:
	\[
	\exists\, t_0 \in (a,b) \colon \quad x_1(t_0) = x_2(t_0).
	\]
	Тогда $x_1(t) \equiv x_2(t)$ на всём $(a,b)$.
\end{Theorem}

\begin{proof}
	Разберём $t \geqslant t_0$; случай $t \leqslant t_0$ симметричен. Положим
	\[
	T = \bigl\{t' \in [t_0, b) \colon\ x_1(t) = x_2(t) \text{ на } [t_0, t']\bigr\}.
	\]
	Множество $T$ непусто, так как $t_0 \in T$, и ограничено сверху числом $b$;
	пусть $t^{*} = \sup T$.

	\medskip
	Если $t^{*} = b$, то решения совпадают на всём $[t_0, b)$ и доказывать
	нечего. Пусть $t^{*} < b$. По непрерывности $x_1$ и $x_2$ из равенства на
	$[t_0, t^{*})$ следует $x_1(t^{*}) = x_2(t^{*})$.

	\medskip
	Точка $t^{*}$ внутренняя для $(a,b)$, поэтому
	$\bigl(t^{*}, x_1(t^{*})\bigr) \in G$ --- точка единственности. Значит,
	найдётся $t' > t^{*}$, для которого $x_1(t) = x_2(t)$ на $[t^{*}, t')$, а
	вместе с равенством на $[t_0, t^{*}]$ это даёт $t' \in T$ --- противоречие с
	тем, что $t^{*} = \sup T$.
\end{proof}

\subsubsection*{Скалярный случай}

Для уравнения $y' = f(x,y)$ условие
\[
f,\ \pder{f}{y} \in C(G)
\]
влечёт $f \in \mathrm{Lip}_{y,\mathrm{loc}}(G)$, то есть выполнены условия
теоремы Пикара. Так обосновываются теоремы существования и единственности,
которыми мы пользовались без доказательства в первой лекции.

\subsection{Метод сжимающих отображений}\label{sec:сжимающие}

Существование решения можно получить и вторым способом --- как неподвижную
точку подходящего оператора.

\begin{Theorem}{Принцип сжимающих отображений (Банах)}{банах}
	Пусть $(X, \rho)$ --- полное метрическое пространство, а $T \colon X \to X$
	--- \textbf{сжимающий оператор}:
	\[
	\exists\, \lambda \in (0,1) \colon \quad
	\rho\bigl(T(x), T(x')\bigr) \leqslant \lambda\,\rho(x,x')
	\qquad \forall x, x' \in X .
	\]
	Тогда у $T$ существует единственная неподвижная точка.
\end{Theorem}

Пусть, как и выше, $\dot x = f(t,x)$, где $f \in C(G)$ и
$f \in \mathrm{Lip}_{x,\mathrm{loc}}(G)$, $(t_0,x_0) \in G$, а
$R = \bigl\{\abs{t - t_0} \leqslant \alpha,\ \abs{x - x_0} \leqslant \beta\bigr\}
\subset G$, $\abs{f} \leqslant M$ в $R$, $L$ --- константа Липшица в $R$ и
$h = \min(\alpha, \beta/M)$. Возьмём дополнительно число $h_0$ так, что
\[
h_0 < h, \qquad L h_0 < 1,
\]
и положим $\Pi = [t_0 - h_0,\ t_0 + h_0]$.

\begin{Statement}{Существование на $\Pi$}{сжатие-существование}
	Существует решение задачи Коши с начальными данными $(t_0,x_0)$
	на промежутке $\Pi$.
\end{Statement}

\begin{proof}
	\textbf{Шаг 1 (пространство).} Положим
	\[
	X = \bigl\{\vphi \colon \Pi \to \RR^n \ \big|\ \vphi \in C(\Pi),\
	(t, \vphi(t)) \in R \text{ при } t \in \Pi\bigr\},
	\qquad
	\rho(\vphi, \vphi') = \max_{t \in \Pi} \abs{\vphi(t) - \vphi'(t)} .
	\]
	Стандартное утверждение анализа: $(X, \rho)$ --- полное метрическое
	пространство.

	\medskip
	\textbf{Шаг 2 (оператор).} Для $\vphi \in X$ положим $\psi = T(\vphi)$, где
	\[
	\psi(t) = x_0 + \int_{t_0}^{t} f\bigl(s, \vphi(s)\bigr)\,\dd s .
	\]
	Оператор действительно переводит $X$ в себя: непрерывность $\vphi$ даёт
	непрерывность $\psi$, а
	\[
	\abs{\psi(t) - x_0} \leqslant M\abs{t - t_0}
	\leqslant M h_0 \leqslant M h \leqslant \beta
	\qquad\text{на } \Pi .
	\]

	\medskip
	\textbf{Шаг 3 (сжатие).} Для $\vphi, \vphi' \in X$
	\[
	\rho\bigl(T(\vphi), T(\vphi')\bigr)
	= \max_{t \in \Pi} \abs{\int_{t_0}^{t}
	\bigl[f\bigl(s, \vphi(s)\bigr) - f\bigl(s, \vphi'(s)\bigr)\bigr]\,\dd s}
	\leqslant L h_0 \max_{t \in \Pi} \abs{\vphi(t) - \vphi'(t)}
	= L h_0\,\rho(\vphi, \vphi') ,
	\]
	а $L h_0 < 1$ по выбору $h_0$. По принципу сжимающих отображений у $T$ есть
	единственная неподвижная точка $\vphi$, а равенство $\vphi = T(\vphi)$ ---
	это и есть интегральное уравнение~\eqref{eq:эиу}.
\end{proof}

\begin{Remark}[addon]{Чем этот способ отличается}{сжатие-сравнение}
	Сравнение с последовательными приближениями поучительно: итерации
	$\vphi_{k+1} = T(\vphi_k)$ --- это в точности итерации того же оператора, а
	лемма о сходимости приближений --- «ручное» доказательство сходимости этих
	итераций. Плата за ссылку на готовый принцип --- лишнее условие
	$L h_0 < 1$: промежуток приходится укорачивать, тогда как оценка
	из~\eqref{eq:оценка-psi} работает на всём промежутке Пеано.
\end{Remark}
